\chapter*{Introduction}
\addcontentsline{toc}{chapter}{Introduction}
\markboth{Introduction}{Introduction}

\section*{Topicality of the Problem}
\addcontentsline{toc}{section}{Topicality of the Problem}

The global automotive diagnostics market was valued at approximately 43~billion USD in~2023 and is forecast to grow at a compound annual rate of~8.2\,\% through~2030~\cite{markets2023}.
A key catalyst is the proliferation of low-cost Bluetooth OBD-II adapters based on the ELM327 microcontroller, which can transform any modern smartphone into a diagnostic instrument for less than 30~USD.
Despite this democratisation of hardware access, the software interpretation layer has not advanced proportionally.

When an Electronic Control Unit (ECU) detects a fault condition, it stores a standardised Diagnostic Trouble Code (DTC) as defined by SAE~J1979 and activates the malfunction indicator lamp.
Retrieving the code via an ELM327 adapter is technically straightforward.
Accurately interpreting it is not.
A single code such as P0171 (\textit{Fuel System Too Lean, Bank~1}) may be caused by a vacuum leak in the intake manifold, a contaminated Mass Airflow (MAF) sensor, insufficient fuel pump pressure, or a failed upstream oxygen sensor---each requiring a fundamentally different repair strategy with costs ranging from negligible to several hundred euros.

Two categories of solutions currently dominate the market.
Traditional OBD readers (e.g., Torque~Pro, Car Scanner~ELM OBD2, OBD Auto~Doctor) successfully retrieve and display fault codes but provide only generic textual descriptions sourced from lookup tables.
They perform no analysis of which specific cause applies to the vehicle under current operating conditions.
AI-augmented tools represent a newer category that addresses the explanation problem; however, they introduce a different failure mode: without access to live sensor readings, large language models (LLMs) produce responses based on their training-set statistics rather than measured vehicle evidence.
A controlled experiment conducted during this research confirmed a cause identification accuracy of only 61\,\% and result consistency of 58\,\% for the direct AI-only approach across three DTC scenarios~\cite{zhang2024gpt}.

This thesis addresses the identified gap by proposing and implementing a \textbf{two-stage diagnostic pipeline}: a deterministic Rule-Based engine verifies the fault cause against live OBD-II sensor data in Stage~1, and generative AI is invoked exclusively in Stage~2 to explain the pre-verified result.

\section*{Aim of the Study}
\addcontentsline{toc}{section}{Aim of the Study}

The aim of this thesis is to design, implement, and experimentally validate a cross-platform mobile diagnostic application---\textit{AI Auto Doctor}---that realises a two-stage analysis pipeline combining deterministic rule-based sensor verification with generative AI explanation, thereby achieving measurably higher diagnostic accuracy, result consistency, and user explainability than existing AI-only or traditional OBD approaches.

\section*{Object and Subject of the Study}
\addcontentsline{toc}{section}{Object and Subject of the Study}

\textbf{Object of the study:} Cross-platform mobile application development for automotive vehicle diagnostics using OBD-II sensor data, rule-based diagnostic algorithms, and generative AI.

\textbf{Subject of the study:} A two-stage diagnostic pipeline combining a deterministic Rule-Based \texttt{DiagnosticEngine} with Google Gemini~1.5~Flash AI, implemented as a Flutter mobile application with BLE OBD-II connectivity, real-time LiveData processing, and structured JSON response format.

\section*{Research Problem and Hypothesis}
\addcontentsline{toc}{section}{Research Problem and Hypothesis}

\textbf{Research problem:} No widely available mobile diagnostic solution implements a structured verification layer that cross-validates DTC codes against live sensor readings before engaging AI explanation, leaving drivers without reliable, actionable, sensor-grounded diagnostic information.

\textbf{Hypothesis:}
\begin{quote}
\itshape
If a rule-based engine first verifies the fault cause against live OBD-II sensor data---including STFT, LTFT, MAF, O2 sensor voltage, coolant temperature, and FreezeFrame snapshot---and generative AI is invoked only with the pre-verified \texttt{DiagnosticResult} as context rather than a raw DTC code, then the resulting two-stage pipeline will achieve statistically significant improvements in cause identification accuracy, result consistency, and user-perceived explainability compared to direct AI-only output.
\end{quote}

\section*{Tasks of the Study}
\addcontentsline{toc}{section}{Tasks of the Study}

The following tasks are defined to achieve the stated aim:

\begin{enumerate}
  \item Conduct a systematic literature review of the OBD-II standard, ELM327 protocol,
        existing mobile diagnostic applications, and AI-based diagnostic approaches
        to identify the research gap (Chapter~\ref{ch:litreview}).
  \item Design the three-layer system architecture (Presentation / Domain / Data) and
        specify the \texttt{DiagnosticEngine} rule sets, \texttt{HealthScoreCalculator},
        BLE communication protocol, SQLite data model, and Gemini AI integration
        (Chapter~\ref{ch:architecture}).
  \item Implement all functional capabilities across eight application screens and
        conduct functional, comparative, and performance testing
        (Chapter~\ref{ch:implementation}).
  \item Verify the research hypothesis through a controlled comparative experiment
        measuring cause accuracy, result consistency, and user explainability for
        AI-Only vs.\ Rule+AI approaches (Chapter~\ref{ch:implementation}).
  \item Formulate conclusions that directly address each task and assess the practical
        and scientific significance of the results.
\end{enumerate}

\begin{table}[!htbp]
\centering
\caption{Mapping of research tasks to thesis chapters}
\label{tab:tasks}
\rowcolors{2}{tblalt}{white}
\begin{tabularx}{0.95\linewidth}{>{\bfseries}c p{9cm} X}
\toprule
\rowcolor{tblhdr}
\color{white}\textbf{Task} &
\color{white}\textbf{Description} &
\color{white}\textbf{Chapter} \\
\midrule
1 & Literature review and research gap identification & Chapter~1 \\
2 & System architecture and component design         & Chapter~2 \\
3 & Implementation and multi-level testing           & Chapter~3 \\
4 & Comparative experiment and hypothesis verification& Chapter~3 \\
5 & Conclusions and significance assessment          & Conclusions \\
\bottomrule
\end{tabularx}
\end{table}

\section*{Research Methods}
\addcontentsline{toc}{section}{Research Methods}

The methodological basis of this thesis comprises the following approaches:

\begin{itemize}
  \item \textbf{Systematic literature review} of automotive diagnostics standards
        (SAE~J1979~\cite{sae2020}, ISO~15765-4~\cite{iso2021}), ELM327 AT command
        protocol~\cite{elm2019}, mobile development frameworks, and AI diagnostic
        literature.
  \item \textbf{Comparative analysis} of twelve existing mobile OBD-II diagnostic
        applications to identify capability gaps.
  \item \textbf{Experimental comparison} of two diagnostic methods on three DTC
        scenarios (P0171, P0300, P0420), each evaluated in three repeated runs by five
        automotive domain experts.
  \item \textbf{Iterative Agile development} using a demo simulation environment
        (\texttt{DemoObdSource}) replicating a Volkswagen Passat~2019~1.4~TSI sensor
        profile.
  \item \textbf{User Acceptance Testing (UAT)} of ten functional requirements on
        Android~14 (Samsung Galaxy~A54) and iOS~17 (Apple iPhone~13).
\end{itemize}

\section*{Thesis Structure}
\addcontentsline{toc}{section}{Thesis Structure}

The thesis is structured as follows.
Chapter~1 provides the theoretical background: OBD-II standard, ELM327 protocol, DTC code structure, analysis of existing diagnostic applications, review of AI and Rule-Based diagnostic approaches, and justification of the Flutter framework.
Chapter~2 presents the system architecture and detailed technical design of all system components.
Chapter~3 describes the implementation of functional capabilities and reports results from functional testing, the comparative hypothesis validation experiment, performance benchmarking, and a STRIDE security analysis.
The Conclusions section addresses each research task and assesses scientific and practical significance.

\clearpage
